\documentclass{article}%
\usepackage[T1]{fontenc}%
\usepackage[utf8]{inputenc}%
\usepackage{lmodern}%
\usepackage{textcomp}%
\usepackage{lastpage}%
\usepackage{geometry}%
\geometry{margin=2cm}%
\usepackage{expex}%
\usepackage[T1]{fontenc}%
\usepackage[utf8]{inputenc}%
\usepackage[spanish]{babel}%
%
\title{Análisis Lingüístico: Gramatica Normativa Mam}%
\author{Generado por el TFM de Elías Carrasco}%
%
\begin{document}%
\normalsize%
\maketitle%
\section{Page 230}%
\label{sec:Page230}%
... Poq'b'il t{-}xilen yol / Extensiones semánticas %
Los cambios semánticos se toman cuando una palabra ha to{-} %
mado otros significados distintos a los originales o comunes. %
Algunas palabras son similares en pronunciación pero vienen %
de diferentes fuentes. %
Las razones más comunes del cambio semántico provienen de %
la pronunciación y/o en extensiones de variantes específicos %
en el contexto que se habla. %
Los cambios semánticos se dan sobre todo en sustantivos, ad{-} %
jetivos y verbos. %
... /0 B'ib'aj/Sustantivos %
Los cambios ocurren en las formas absolutas de las palabras %
como en su forma poseída, la mayoría de los cambios se deben %
a procesos fonológicos. %
Ejemplos: 

%
\end{document}